\honest{The Parable of Hemlock}{Eliezer Yudkowsky}{2008}

I open with ``All men are mortal. Socrates is a man. Therefore Socrates is mortal,'' attributed to ``Aristotle(?)''. \nb{The doubt is warranted. The example is not in Aristotle; the earliest printing David Wheeler could find is in John Stuart Mill's \textsc{A System of Logic} (1843).}

A parable. As Socrates drinks the hemlock, a student of philosophy assures an onlooker that he will die, because all men are mortal ``by definition.'' They wait, and Socrates does not die. ``Then Socrates must not be a man,'' says the student, and calls this too ``a logical certainty.''

My thesis: ``you can't make reality go a different way by choosing a different definition.'' You could reason from experience instead: the last thirty humans who drank hemlock fell over, Socrates looks like a human, so he will fall over too. But that is a guess, and the Greek philosophers were ``rather fond of certainty.'' So they reply that mortality is part of the definition of ``human.'' \nb{Porphyry, in the third century AD, did define man as a rational mortal animal, and medieval logicians adopted the definition.} Then you cannot know that Socrates is human until you see him die: fluent Greek, red blood, even human DNA are not logically equivalent to mortality. Even then you cannot be certain. He might rise from the grave, or be signed up for cryonics. If mortality means a finite life, you must watch to the end of eternity. Or you might have been deceived by a projected illusion, or hallucinated it. \nb{This dilemma is Mill's, in the same chapter and with the same syllogism: every syllogism, taken as proof, begs the question, and ``All inference is from particulars to particulars.'' The post does not mention Mill, whose argument supports it.}

The problem with syllogisms is that ``they're always valid.'' A valid syllogism is valid in every possible world, including neighboring worlds where hemlock is a delicious treat and worlds without Socrates. Evidence for a hypothesis is what we are more likely to see if it is true than if it is false, so a syllogism's validity is no evidence about which world we live in.

Logic is not useless. It tells us what, in a sense, we already know, and we do not always believe what we know. Whether 29384209 is prime is settled by my own axioms, but I must still work it out. Likewise, from the uncertain guesses that humans are vulnerable to hemlock and that Socrates is human, logic tells me that my guesses predict Socrates is vulnerable. \nb{It is not prime: 1667 times 17627.} ``It's been suggested'' that logical reasoning removes our uncertainty about ``impossible possible worlds.'' \nb{The phrase comes from the title of a 1975 paper by Jaakko Hintikka, ``Impossible Possible Worlds Vindicated.''} In that sense logic works like observation. But whether Socrates will keel over, or do fifty jumping jacks and compete in the Olympics, is a question about possible worlds, not impossible ones. Logic can bring old observations and guesses to bear and say what they predict, but it never settles by itself a real-world question that could go either way. Last, a ``Litany Against Logic'': ``Logic stays true, wherever you may go, / So logic never tells you where you live.''
